\section{Podsumowanie stanu wiedzy i wnioski dla projektu}
\label{subsubsec:llm_acr_synthesis}

Najnowsze publikacje kreują obraz tego w jakim obecnie punkcie znajduje się zagadnienie automatycznego
przeglądu kodu. Widać głównie, że przestano traktować LLM jako cudowne i samodzielne narzędzie, które samo
ma rozwiązać problemy z kodem. Zamiast tego, nowoczesne narzędzia idą w stronę złożonych architektur systemowych, co pokazują 
Meng i Ramesh \cite{Meng2025RARe, Ramesh2025EricssonExperience}. Model językowy staje się jednym z wielu modułów w całym
ekosystemie.

Do stworzenia właśnego systemu przeglądu kodu, można wyciągnąć parę konkretnych ustaleń:
\begin{itemize}
    \item \textbf{Samo promptowanie to za mało.} Choć techniki opisane w rozdziale \ref{subsec:llm_acr_prompting} 
są to najszybsze i najprostsze rozwiązanie, to ich ogromna wrażliwość 
na drobne zmiany w treści zapytania jest problematyczna. Zamiast statycznych instrukcji 
lepiej postawić na ustrukturyzowane ICL, gdzie kontekst dobierany jest automatycznie (np. przez BM25), co 
sugerują Pornprasit i Haider \cite{Pornprasit2024FineTuningPromptingCR, Haider2024PromptingFineTuningRCG}.
    
    \item \textbf{Znaczenie jakości danych nad architekturą.} Jak już wspomniano przy omawianiu fine-tuningu 
(\ref{subsec:llm_acr_finetuning}), metody takie jak LoRA czy QLoRA faktycznie pomagają "wymusić" na modelu konkretny styl 
czy strukturę odpowiedzi. Kluczowe nie jest tylko dostrajanie parametrów, a raczej to, co 
do modelu dodajemy. Metoda DRC pokazuje, jak ważne jest odfiltrowanie szumu z "ludzkich" zbiorów 
danych, żeby model nie uczył się złych nawyków \cite{Yu2025Desiview, Yu2024Carllm}.
    
    \item \textbf{Kontekst repozytorium jest bardzo ważny.} Patrzenie tylko na 
same \emph{diffy} to zły kierunek, bo model ma wtedy zbyt wąski horyzont wiedzy \cite{Tufano2024CodeReviewAutomation}. 
Skuteczny system ACR musi wiedzieć więcej, czyli integrować RAG z analizą statyczną. 
To właściwie jedyny sposób, żeby ograniczyć generowanie halucynacji i zbyt ogólnego generowania tekstu w komentarzach \cite{Jaoua2025LLMStaticAnalyzersCR, 
Icoz2025SymbolicReasoning}. Trochę więcej miejsca poświęcono temu w sekcji \ref{subsec:llm_acr_context_hybrid}.
    
    \item \textbf{Zmiana myślenia o ewaluacji.} Tradycyjne metryki, którymi często ocenia się tekst 
(jak BLEU czy BERTScore), w przypadku kodu słabo się sprawdzają, o czym była mowa 
w \ref{subsec:llm_acr_evaluation_gaps}. W projektowaniu systemu priorytetem powinna być weryfikacja funkcjonalna i to, czy 
sugerowane zmiany faktycznie przechodzą testy jednostkowe. Chodzi o to, żeby system był faktycznie użyteczny \cite{Cihan2025EvaluatingLLMsForCodeReview, Lin2025CodeReviewQA}.
    
    \item \textbf{Problem wielowymiarowości i jakości komentarzy.} Code review nie jest
zadaniem jednowymiarowym - badania pokazują, że 78,7\% komentarzy dotyczy co najmniej
dwóch różnych kategorii problemów \cite{Li2025CodeDoctor}. Skuteczny system musi więc
oceniać komentarze pod kątem trafności (\emph{relevance}), użyteczności
(\emph{actionability}) oraz zrozumiałości (\emph{comprehensibility}). Problem ogólnikowości
(\emph{vagueness}) dotyka aż 61\% komentarzy generowanych tradycyjnymi metodami
\cite{Ren2025HydraReviewer}, przez co w praktyce często niewiele użytecznego z nich wynika.

    \item \textbf{Ryzyko regresji i konieczność warstwy zabezpieczeń.} Nawet zaawansowane
modele potrafią zaproponować zmiany pogarszające poprawny kod - \emph{Regression Ratio}
wynosi 10-13\% \cite{Cihan2025EvaluatingLLMsForCodeReview}. Dodatkowo lokalizacja
fragmentu kodu okazuje się nproblemem modeli LLM \cite{Lin2025CodeReviewQA}.
To wymusza podejście \emph{human-in-the-loop} oraz mechanizmy walidacji funkcjonalnej.
\end{itemize}

Współczesne badania \cite{Tang2025MoEReviewer, Lu2023LLaMAReviewer} coraz częściej odchodzą od traktowania predykcji, generowania i poprawiania jako 
całkowicie osobnych kroków. Architektury wielozadaniowe (np. MoE) są pod tym względem preferowane - łączą 
one te elementyw jednym podejściu. W praktyce ma to sens zwłaszcza wtedy, gdy taki 
system ma działać w potokach CI/CD: da się utrzymać dobrą efektywność, a jednocześnie ograniczać 
koszty.

\paragraph{Główne problemy i braki w obecnych badaniach.}
Widać ogromny postęp w literaturze, jednak przyglądając się publikacjom bliżej, można dostrzec 
kilka dość istotnych luk. Widać sporą dysproporcję w tym, czym zajmują się 
badacze - aż 91,3\% prac skupia się tylko na samym wykrywaniu podatności, podczas gdy 
automatyczna naprawa kodu czy chociażby sensowne wyjaśnianie błędów to wciąż tematy rzadko poruszane \cite{Germano2025VulnSLR}. 
Inna sprawa to kwestia rzetelności samych danych. Okazuje się, że popularne zbiory treningowe są 
mocno "zaśmiecone" (szum sięga tu ok. 25\%), a problem \emph{data leakage} potrafi sztucznie zawyżyć 
wyniki modeli nawet o ponad 30 punktów procentowych \cite{Tufano2024CodeReviewAutomation, Germano2025VulnSLR}. Dochodzą bariery 
czysto praktyczne, o których wspomina się rzadziej, a które w rzeczywistym środowisku programistycznym są 
kluczowe. Chodzi tu głównie o koszty utrzymania takich modeli, czas oczekiwania na odpowiedź i 
to, jak w ogóle wpiąć takie narzędzie w codzienny workflow programisty, żeby nie przeszkadzało 
w pracy \cite{Ramesh2025EricssonExperience}.

\paragraph{Analiza porównawcza podejść - co działa, a co nie.}
Na podstawie przeglądu literatury można wyciągnąć konkretne wnioski dotyczące efektywności poszczególnych
podejść w automatycznym code review. Tabela \ref{tab:porownanie_podejsc} syntetycznie zestawia kluczowe
charakterystyki i wskazuje, kiedy dane rozwiązanie jest optymalne, a kiedy zawodzi.

\begin{table}[h]
\centering
\caption{Porównanie podejść do automatycznego code review z LLM - synteza wniosków z literatury.}
\label{tab:porownanie_podejsc}
\small
\begin{tabular}{p{2.4cm}p{3.8cm}p{3.8cm}p{3.8cm}}
\hline
\textbf{Podejście} & \textbf{Co działa najlepiej} & \textbf{Ograniczenia} & \textbf{Warunki sukcesu} \\
\hline
\textbf{Promptowanie (ICL)} & 
Szybkie prototypowanie; BM25 do doboru przykładów; top-1 few-shot \cite{Pornprasit2024FineTuningPromptingCR, Meng2025RARe} & 
Wrażliwość na konstrukcję promptu; ograniczony kontekst; niska stabilność \cite{Haider2024PromptingFineTuningRCG} & 
Ustrukturyzowane ICL z automatycznym doborem przykładów; nie więcej niż 3 demonstracje \\
\hline
\textbf{Fine-tuning (PEFT)} & 
Spójność stylu komentarzy; efektywność zasobowa (LoRA \textasciitilde16MB) \cite{Lu2023LLaMAReviewer}; ukierunkowanie domenowe & 
Zależność od jakości danych (szum \textasciitilde25\%); ryzyko przeuczenia \cite{Tufano2024CodeReviewAutomation, Yu2025Desiview} & 
Destylacja DRC; filtracja danych; wielozadaniowe architektury (MoE) \cite{Tang2025MoEReviewer, Yu2025Desiview} \\
\hline
\textbf{RAG} & 
Redukcja ogólnikowości (61\% → 16.4\%); wzrost BLEU +110-153\% \cite{Ren2025HydraReviewer, Meng2025RARe} & 
Top-5 gorsze niż top-1 (szum w kontekście); koszt retrievalu; jakość indeksu \cite{Meng2025RARe} & 
Top-1 retrieval; selektywny dobór kontekstu (AST, historia, standardy); deduplikacja \\
\hline
\textbf{Hybryda (LLM+statyka)} & 
Łączenie pokrycia LLM z precyzją reguł; weryfikowalność sugestii \cite{Jaoua2025LLMStaticAnalyzersCR, Ramesh2025EricssonExperience} & 
Złożoność integracji; ryzyko konfliktu sygnałów; koszt utrzymania pipeline & 
RAG oparty o wyniki analizatorów; walidacja funkcjonalna (testy); human-in-the-loop \\
\hline
\end{tabular}
\end{table}

Kluczowym wnioskiem z tego zestawienia jest to, że \textbf{nie ma jednego uniwersalnego podejścia}.
Efektywność zależy od kontekstu wdrożenia: promptowanie działa dobrze w scenariuszach cold-start
i prototypowania, ale stabilność wymaga fine-tuningu na wysokiej jakości danych. RAG radykalnie
poprawia trafność, lecz wymaga przemyślanej strategii retrievalu (top-1 lepsze niż top-k). Podejścia
hybrydowe osiągają najlepsze wyniki pod warunkiem dobrej integracji z workflow, co potwierdza
doświadczenie Ericsson \cite{Ramesh2025EricssonExperience}.

Z perspektywy projektowania systemu decydujące okazują się czynniki praktyczne: koszt inferencji
(MoE: 16.38 TFLOPs vs LLaMA 7B: 57.34 TFLOPs \cite{Tang2025MoEReviewer}), czas odpowiedzi
(kluczowy w CI/CD), oraz możliwość weryfikacji wyników (Regression Ratio 10-13\%
\cite{Cihan2025EvaluatingLLMsForCodeReview}). Literatura jednoznacznie wskazuje, że metryki tekstowe
(BLEU) nie korelują z użytecznością w praktyce - 21.83\% odrzuconych przez EM predykcji było
semantycznie poprawnych \cite{Tufano2024CodeReviewAutomation}. Zamiast tego priorytet powinny mieć
kryteria jakościowe: trafność lokalizacji (\emph{CL}), użyteczność (\emph{actionability}) oraz
ryzyko regresji, co pokazują benchmarki typu CodeReviewQA \cite{Lin2025CodeReviewQA}.

\paragraph{Wnioski i kierunek realizacji projektu.}
Z analizy źródeł wynika, że jeśli system wspomagający code review ma faktycznie działać, nie 
może opierać się na samym modelu LLM. Potrzebna jest tu raczej architektura hybrydowa. Model 
powinien pełnić rolę generatora, ale musi być obudowany dodatkowymi mechanizmami, takimi jak RAG czy 
analiza statyczna i walidacja funkcjonalna.

\paragraph{Identyfikacja luki badawczej.}
Przegląd literatury pokazuje lukę o charakterze praktycznym i systemowym. Istniejące prace koncentrują się na 
pojedynczych elementach: generacji komentarzy \cite{Pornprasit2024FineTuningPromptingCR, Haider2024PromptingFineTuningRCG},
dostrajaniu modeli \cite{Lu2023LLaMAReviewer, Tang2025MoEReviewer} lub RAG \cite{Meng2025RARe}. Brakuje 
rozwiązań end-to-end łączących retrieval kontekstu, generację LLM i walidację funkcjonalną w realnym workflow PR/MR 
z konfigurowalnością per-projekt i warstwą human-in-the-loop. Wdrożenia przemysłowe (Ericsson \cite{Ramesh2025EricssonExperience}) nie adresują wieloplatformowości, prace 
akademickie (RARe \cite{Meng2025RARe}, Hydra \cite{Ren2025HydraReviewer}) nie są osadzone w produkcyjnym GitHub/GitLab lub podobnym.

W ramach pracy luka będzie wypełniana przez system ACR integrujący się z PR/MR jako 
krok CI, wykorzystujący RAG (top-1 retrieval \cite{Meng2025RARe}), łączący LLM z walidacją funkcjonalną (analiza statyczna 
\cite{Jaoua2025LLMStaticAnalyzersCR}, human-in-the-loop \cite{bacchelli2013}), konfigurowalny per-projekt i ewaluowany
kryteriami jakościowymi \cite{Lin2025CodeReviewQA, Ramesh2025EricssonExperience}.

\paragraph{Konkrety techniczne z literatury.}
Przegląd prac pokazuje szereg zaleceń odnośnie rozwiązań technicznych. Warstwa LLM jest konfigurowalna — w projektowanym systemie 
wykorzystane będą gotowe, mocne modele dostępne przez API \cite{Pornprasit2024FineTuningPromptingCR, Lin2025CodeReviewQA}. 
Implementacje RAG stosują retrieval top-1 do top-5 z użyciem baz wektorowych (FAISS) \cite{Meng2025RARe} 
lub retrieval leksykalnego (BM25) \cite{Pornprasit2024FineTuningPromptingCR}. Ekstrakcja kontekstu kodu 
opiera się na Tree-sitter (parser AST) do budowy call graph i izolacji funkcji \cite{Ramesh2025EricssonExperience, Pornprasit2024FineTuningPromptingCR, Ren2025HydraReviewer}. 
Ewaluacja bazuje na BLEU-4 (typowo 5.58--13.13 dla SOTA) \cite{Pornprasit2024FineTuningPromptingCR, Meng2025RARe}, 
BERTScore oraz ocenie praktyków w skali 1--5 (relevance, information, clarity) \cite{Pornprasit2024FineTuningPromptingCR, Tang2025MoEReviewer}. 
Integracja z platformami VCS realizowana przez REST API + webhooks \cite{Ramesh2025EricssonExperience, Jaoua2025LLMStaticAnalyzersCR}, 
co umożliwia obsługę GitHub, GitLab i innych systemów PR/MR. Szczegóły architektury i metodyki ewaluacji opisano w kolejnym rozdziale.
